\section{Módulo 5: Bodegas, Centros de Distribución y Decisiones Estratégicas}

\subsection{Logística y centros de distribución}

\begin{definicion}{Logística y el rol del CD}
\textbf{Logística:} tener el producto/servicio preciso, en el lugar preciso, en el momento preciso, al costo mínimo. El \textbf{tiempo para completar una orden} tiene tres ventanas: recepción de órdenes, procesamiento y transporte. El gran \textbf{trade-off} de bodegas y CD es \textbf{Tiempo vs.\ Espacio}.
\end{definicion}

\begin{definicion}{Beneficios de un Centro de Distribución}
\begin{itemize}[leftmargin=1.4em, itemsep=1pt]
  \item \textbf{Consolidación geográfica:} reduce costos de transporte (menos arcos: con centro \$10 vs.\ sin centro \$25).
  \item \textbf{Consolidación temporal} de la demanda.
  \item \textbf{Agregar valor / servicios:} ensamblaje, \textit{postponement} (DELL), etiquetado (vinos), kitting, control aduanero.
  \item \textbf{Logística reversa} (Ley REP en Chile).
\end{itemize}
Costos: infraestructura y personal. Decisiones: externalizar (construir/dueño vs.\ \textbf{3PL}), localización, tipo (\textit{cross-dock} vs.\ bodega), tamaño.
\end{definicion}

\subsection{Postponement y Risk Pooling}

\begin{definicion}{Postponement (postergación)}
Retrasar la diferenciación del producto hasta el último momento posible reduce el inventario. \textbf{Ejemplo:} 4 procesadores $\times$ 4 memorias $\times$ 4 tarjetas = 64 PC terminados $\Rightarrow$ con 10 c/u son 640 en inventario. Si \textbf{postpongo} el ensamble: $4\times10 + 4\times10 + 4\times10 = 120$ componentes en inventario con el \textbf{mismo} nivel de servicio.
\end{definicion}

\begin{definicion}{Risk Pooling (agregación de riesgo)}
Centralizar el inventario de varios mercados reduce la variabilidad relativa (CV) de la demanda agregada, bajando el stock de seguridad total. La reducción es mayor mientras \textbf{menos correlacionadas} estén las demandas. Conecta con \textbf{Newsvendor / SS} de la I1.
\end{definicion}

\subsection{Análisis de flujo y localización}

\begin{formulas}{Ley de flujo en bodegas ($Q = A\cdot v$)}
\[
  Q = A \times v
\]
$Q$ = caudal/volumen (capacidad de mover inventario, rotación), $A$ = área/tamaño de bodega en unidades, $v$ = velocidad de flujo o rotación (unidades/tiempo). Sirve para dimensionar bodegas: dada la demanda anual y la capacidad, se despeja la rotación $v$ y de ahí el personal/montacargas requerido.
\end{formulas}

\begin{formulas}{Localización: modelo de distancia euclidiana}
Minimizar la suma ponderada de distancias a instalaciones existentes:
\[
  \min f(X) = \sum_{i=1}^{m} w_i\, d(X, P_i), \qquad
  d(X,P_i) = 111\sqrt{(x-a_i)^2 + (y-b_i)^2}\ \text{[km, con lat/long]}
\]
Factor $111$ km $\approx 69$ millas por grado. Para optimizar se usa la versión cuadrática (centro de gravedad). También existe distancia \textbf{Manhattan} (rectilínea).
\end{formulas}

\subsection{V/F frecuentes de Bodegas/CD}

\begin{center}
\small
\begin{tabular}{@{}clp{10.5cm}@{}}
  \toprule
  $\#$ & R & Afirmación y corrección \\
  \midrule
  1 & \textbf{V} & ``El postponement reduce el inventario manteniendo el mismo nivel de servicio.'' \\
  2 & \textbf{F} & ``El risk pooling funciona mejor cuando las demandas están muy correlacionadas.'' --- Falso: funciona mejor con demandas \textbf{poco correlacionadas / independientes}. \\
  3 & \textbf{V} & ``Un centro de distribución reduce el costo de transporte por consolidación geográfica.'' \\
  4 & \textbf{F} & ``Un cross-dock almacena inventario por largos períodos.'' --- Falso: el cross-dock \textbf{minimiza} el almacenaje (recibe y despacha rápido). \\
  \bottomrule
\end{tabular}
\end{center}

% ════════════════════════════════════════════════════════════════════════════



% ════════════════════════════════════════════════════════════════════════════
